\section{Introdução e Contextualização} \label{section: introduction}

PREENCHER: contexto biológico do alvo e sua relevância terapêutica; anatomia do sítio de ligação (sub-sítios e resíduos, de preferência com uma figura didática); origem da série estudada e os dados experimentais que motivam a modelagem; a pergunta central do estudo.

As estruturas das moléculas discutidas ao longo do relatório estão reunidas na \cref{tab:estruturas}, que serve de guia para a leitura das seções seguintes.

\begin{table}[p]
\centering
\caption{Estruturas 2D das moléculas estudadas, geradas a partir dos SMILES usados no \textit{docking}. PREENCHER: modificação estrutural em relação ao composto líder e atividade medida.}
\label{tab:estruturas}
\renewcommand{\arraystretch}{1.1}
\begin{tabular}{>{\centering\arraybackslash}m{0.45\textwidth} >{\centering\arraybackslash}m{0.45\textwidth}}
PREENCHER: figura da pasta Figures/Estruturas \par (a) nome --- modificação --- atividade &
PREENCHER \par (b) nome --- modificação --- atividade \\
\end{tabular}
\end{table}

PREENCHER: tabela com os dados experimentais (atividade, concentração, controles do ensaio).
